\label{sec:montana}

\subsection{Background and Procedural History}

After the 1990 Census, Montana's population (799{,}065) entitled it to
$q_{\text{MT}} = 799{,}065 \times 435 / 248{,}709{,}873 \approx 1.398$ seats.
Under Huntington-Hill, Montana's priority for a second seat was
$P_{\text{MT}}(1) = 799{,}065/\sqrt{2} \approx 565{,}034$,
which fell below the priority cutoff for the 435th seat.
Montana received 1 seat; Washington received the last seat.

Montana brought suit arguing that the Huntington-Hill method violated
Article~I, \S2's requirement of proportional apportionment because
a single-member Montana district with 799{,}065 residents was
``as equal as possible'' in population to other districts, and that
a method giving Montana 1 seat while Washington (with comparable
population) received more was arbitrary.
Montana proposed the Dean method (an alternative divisor method)
or a ``maximum deviation'' standard as alternatives.

\subsection{Procedural Posture}

The district court found for Montana, holding that Huntington-Hill
was not the method that minimised the maximum deviation from equal
population.
The United States appealed directly to the Supreme Court under
28~U.S.C.\ \S1253 (direct appeal in reapportionment cases).

\subsection{The Supreme Court's Holding}

The Supreme Court reversed in \textit{United States Dept.\ of Commerce v.\
Montana}, 503~U.S.\ 442 (1992), in a unanimous opinion authored by
Justice Stevens \citep{montana1992}.
The Court held:

\begin{enumerate}
  \item \textbf{Political question doctrine does not bar review}:
    The case presented a justiciable question about statutory
    interpretation, not a purely political question.
  \item \textbf{Congress has broad discretion}: Article~I, \S2 does not
    specify a mathematical method; Congress may choose any method that
    is a reasonable interpretation of ``proportional'' representation.
    The Court rejected Montana's argument that only one method
    (the Dean method or maximum-deviation minimisation) satisfied
    the Constitution.
  \item \textbf{Huntington-Hill is constitutional}: The equal proportions
    method is a rational choice by Congress.
    It has mathematical properties (the relative population pair test)
    that justify it as a form of proportionality.
    The Court noted that the debate among mathematicians about which
    method is ``best'' is one that courts are ill-equipped to resolve,
    and that statutory deference to Congress is appropriate.
  \item \textbf{Fractions make perfection impossible}: The Court
    acknowledged that any apportionment of indivisible seats will
    produce some inequity, and that Congress is not required
    to minimise any particular measure of that inequity as long as
    its chosen method is rational.
\end{enumerate}

\subsection{Significance for \textsc{bisect-apportion}}

\textit{Montana} establishes three points relevant to the
\textsc{bisect-apportion} implementation:

First, the correct implementation of Huntington-Hill is not merely
a mathematical exercise but a legal requirement.
The statutory apportionment method determines how many districts each
state draws; any error in the implementation produces downstream
redistricting errors.

Second, the Court's discussion of alternative methods
(Dean, Adams, Webster) confirms that these methods are mathematically
distinct from Huntington-Hill in ways that can affect specific states.
\textsc{bisect-apportion}'s current implementation of the four divisor methods
enables
the ``what if Congress had chosen Webster?'' counterfactual analysis
that practitioners and courts may request.

Third, \textit{Montana}'s political question discussion confirms that
the mathematical properties of apportionment methods are justiciable
issues.
Future SHA-verified Census fixtures for \textsc{bisect-apportion} will be
legally as well as mathematically significant: they would establish an auditable
chain of custody from the Census Bureau's published population figures to the
final seat counts.
